\chapter{Протоколы с нулевым разглашением}
\label{chapter:zk-protocols}
% Intro from https://core.ac.uk/download/pdf/82198769.pdf
% CONTAINS NICE HIHG LEVEL STUFF
% Term SNARK was introduced in  Bitansky, Nir; Canetti, Ran; Chiesa, Alessandro; Tromer, Eran (January 2012). ``From extractable collision resistance to succinct non-interactive arguments of knowledge, and back again''
Так называемый \term{протокол с нулевым разглашением} (zero-knowledge protocol) --- это набор математических правил, с помощью которых одна сторона, обычно называемая \term{доказывающим} (prover), может убедить другую сторону, обычно называемую \term{проверяющим} (verifier), что при данном экземпляре (instance) доказывающий знает некоторый свидетель (witness) для этого экземпляра, не раскрывая никакой информации о свидетеле.

Как мы видели в \chaptname{} \ref{chap:statements}, для данного языка $L$ и экземпляра $I$ утверждение о знании <<существует свидетель $W$ такой, что $(I;W)$ является словом в $L$>> конструктивно доказывается путём предоставления свидетеля $W$ проверяющему. Проверяющий затем может использовать грамматику языка для проверки доказательства. Напротив, задача любого протокола с нулевым разглашением заключается в том, чтобы позволить доказывающему доказать знание свидетеля любому проверяющему, не раскрывая никакой информации о свидетеле, кроме его существования.

В этой главе мы рассмотрим различные системы, существующие для решения этой задачи. Мы начнём с введения в основные понятия и терминологию систем доказательства с нулевым разглашением, а затем представим так называемый протокол Groth\_16 как одну из наиболее эффективных систем. В будущих версиях этой книги мы обновим эту главу другими системами доказательства с нулевым разглашением.

\section{Системы доказательства}
С абстрактной точки зрения, система доказательства --- это набор правил, моделирующих генерацию и обмен сообщениями между двумя сторонами, обычно называемыми доказывающим и проверяющим. Цель системы доказательства состоит в том, чтобы установить, принадлежит ли данная строка формальному языку или нет.

Системы доказательства часто классифицируются по определённым предположениям о доверии и вычислительным возможностям доказывающего и проверяющего. В наиболее общей форме доказывающий обычно обладает неограниченными вычислительными ресурсами, но ему нельзя доверять, в то время как проверяющий имеет ограниченную вычислительную мощность, но предполагается честным.

Доказательство принадлежности или утверждений о знании для некоторой строки, как объяснялось в \chaptname{} \ref{chap:statements}, выполняется путём генерации определённых сообщений, которые отправляются между доказывающим и проверяющим, пока проверяющий не убедится, что строка является словом в рассматриваемом языке.

% https://link.springer.com/content/pdf/10.1007/BF00195207.pdf
Более конкретно, пусть $\Sigma$ --- алфавит, и пусть $L$ --- формальный язык, определённый над $\Sigma$. Тогда \term{система доказательства} (proof system) для языка $L$ --- это пара вероятностных интерактивных алгоритмов $(P,V)$, где $P$ называется \term{доказывающим} (prover), а $V$ называется \term{проверяющим} (verifier).

Оба алгоритма могут отправлять сообщения друг другу, каждый алгоритм имеет своё собственное состояние, некоторое общее начальное состояние и доступ к сообщениям. Проверяющий ограничен числом шагов, полиномиальным от размера общего начального состояния, после чего он останавливается и выдаёт либо \texttt{accept}, либо \texttt{reject}, указывая, принимает ли он данную строку как слово в $L$ или нет. Напротив, в наиболее общей форме системы доказательства на вычислительную мощность доказывающего не накладывается никаких ограничений.

Когда выполнение алгоритма проверяющего завершается, должны выполняться следующие условия:
\begin{itemize}
\item (Полнота) Если строка $x\in \Sigma^*$ является словом в языке $L$ и обе стороны следуют протоколу, проверяющий выдаёт \texttt{accept}.
\item (Корректность) Если строка $x\in \Sigma^*$ не является словом в языке $L$ и проверяющий следует протоколу, проверяющий выдаёт \texttt{reject}, за исключением некоторой малой вероятности.
\end{itemize}
Кроме того, система доказательства называется \term{с нулевым разглашением} (zero-knowledge), если проверяющий не узнаёт ничего о $x$, кроме того, что $x\in L$.

Предыдущее определение систем доказательства очень общее, и в этой области известно много подклассов систем доказательства. Например, некоторые системы доказательства ограничивают вычислительную мощность доказывающего, в то время как другие предполагают, что проверяющий имеет доступ к случайности. Кроме того, системы доказательства классифицируются по количеству сообщений, которые могут быть обменены. Если системе требуется отправить лишь одно сообщение от доказывающего к проверяющему, система доказательства называется \term{неинтерактивной} (non-interactive), поскольку никакого взаимодействия, кроме отправки самого доказательства, не требуется. Напротив, любая другая система доказательства называется \term{интерактивной} (interactive).

Система доказательства обычно называется \term{краткой} (succinct), если размер доказательства меньше, чем свидетель, необходимый для генерации доказательства. Более того, система доказательства называется \term{вычислительно корректной} (computationally sound), если корректность выполняется лишь при предположении, что вычислительные возможности доказывающего полиномиально ограничены. Чтобы отличать общие доказательства от вычислительно корректных доказательств, последние часто называют \term{аргументами} (arguments).

Термин \term{zk-SNARK} является аббревиатурой от <<Zero-knowledge, succinct, non-interactive argument of knowledge>> (аргумент знания с нулевым разглашением, краткий, неинтерактивный). Следовательно, системы доказательства, способные генерировать zk-SNARK, обладают свойством нулевого разглашения (zero-knowledge), способны генерировать доказательства, требующие меньше места, чем исходный свидетель, и не требуют взаимодействия между доказывающим и проверяющим, кроме передачи самого zk-SNARK. Однако эти системы корректны лишь при предположении, что вычислительные возможности доказывающего полиномиально ограничены.
\begin{example}[Конструктивные доказательства для алгебраических схем] Как мы видели в \ref{circuit-satisfiability}, алгебраические схемы порождают формальные языки и конструктивные доказательства для утверждений о знании.

Чтобы переформулировать это понятие конструктивных доказательств для алгебраических схем в систему доказательства, пусть $\F$ --- конечное поле, и пусть $C(\F)$ --- алгебраическая схема над $\F$ с ассоциированным языком $L_{C(\F)}$. Неинтерактивная система доказательства для $L_{C(\F)}$ задаётся следующими двумя алгоритмами:

\term{Алгоритм доказывающего}: Доказывающий $P$ определяется выполнением схемы. При данном экземпляре $I$ доказывающий выполняет схему $C(\F)$ для вычисления свидетеля $W$ такого, что пара $(I;W)$ является допустимым назначением для $C(\F)$, если схема выполнима для $I$. Доказывающий затем отправляет конструктивное доказательство $(I;W)$ проверяющему.

\term{Алгоритм проверяющего}: Получив сообщение $(I;W)$, алгоритм проверяющего $V$ подставляет $(I;W)$ в ассоциированную R1CS схемы. Если $(I;W)$ является решением R1CS, проверяющий возвращает \texttt{accept}, если нет --- \texttt{reject}.

Чтобы убедиться, что эта система доказательства полна и корректна, пусть $C(\F)$ --- схема над полем $\F$, и пусть $I$ --- экземпляр. Схема может иметь или не иметь свидетеля $W$ такого, что $(I;W)$ является допустимым назначением для $C(\F)$.

Если $W$ не существует, $I$ не является частью никакого слова в $L_{C(\F)}$, и нет способа для $P$ сгенерировать допустимое назначение. Следовательно, проверяющий не примет никакое заявленное доказательство, отправленное $P$, поскольку ассоциированная R1CS не имеет решений для экземпляра $I$. Это означает, что система \term{корректна} (sound).

Если, с другой стороны, $W$ существует и $P$ честен, $P$ может использовать свою неограниченную вычислительную мощность для вычисления $W$ и отправить $(I;W)$ к $V$, который примет его, поскольку это решение ассоциированной R1CS. Это означает, что система \term{полна} (complete).

Система неинтерактивна, поскольку доказывающий отправляет лишь одно сообщение проверяющему, которое содержит само доказательство. Однако система доказательства \term{не} является краткой, поскольку доказательство является свидетелем. Система доказательства также не обладает свойством нулевого разглашения, поскольку проверяющий имеет доступ к свидетелю и, следовательно, узнаёт о нём всё.
\end{example}
\section{Протокол Groth16}
\label{sec:gorth_16}
 В \chaptname{} \ref{chap:statements} мы представили алгебраические схемы, их ассоциированные системы ограничений ранга 1 (Rank-1 Constraint Systems) и порождённые ими квадратичные арифметические программы (Quadratic Arithmetic Programs). Эти модели определяют формальные языки, а ассоциированные утверждения о принадлежности и знании предоставляют конструктивные доказательства путём выполнения схемы для вычисления решения ассоциированной R1CS. Как отмечалось ранее в разделе \ref{sec:QAP}, доказательство затем может быть преобразовано в многочлен, который делится на другой многочлен только в том и только в том случае, если доказательство корректно.
 
В \cite{Groth16} Йенс Грот (Jens Groth) разработал метод преобразования конструктивных доказательств в краткие неинтерактивные аргументы знания с нулевым разглашением (zero-knowledge succinct non-interactive arguments of knowledge). Заданные группы $\G_1$, $\G_2$ и $\G_3$, и эффективно вычислимое отображение спаривания $e(\cdot,\cdot): \G_1 \times \G_2 \to \G_3$ (см. \ref{pairing-map}), результирующий zk-SNARK в протоколе Грота имеет постоянный размер, состоящий из двух элементов из $\G_1$ и одного элемента из $\G_2$, независимо от размера экземпляра и свидетеля. Проверка неинтерактивна и требует вычисления числа возведений в степень, пропорционального размеру экземпляра, вместе с тремя групповыми спариваниями, чтобы проверить одно доказательство.

Сгенерированный zk-SNARK обладает свойством нулевого разглашения, имеет полноту и корректность в общей модели билинейных групп, при предположении существования доверенной третьей стороны, которая выполняет предобработку для создания общей опорной строки (Common Reference String) и симуляционного люка (simulation trapdoor). Крайне важно, чтобы эта сторона была доверенной в удалении симуляционного люка, поскольку любой, кто обладает им, имел бы возможность моделировать доказательства. Как показано в \cite{bowe-17}, можно преобразовать одностороннюю доверенную установку (trusted setup) в многостороннее вычисление, которое является безопасным, пока хотя бы один участник удаляет свой вклад в симуляционный люк.

Более точно, пусть $R$ --- система ограничений ранга 1 (Rank-1 Constraint System), определённая над некоторым конечным полем $\F_r$. Тогда \term{параметры Groth16} (Groth16 parameters) для $R$ задаются следующим множеством:
\begin{equation}
\label{groth16-parameters}
\mathtt{Groth16\text{-}Param}(R)= \{ r, \G_1, \G_2, e(\cdot,\cdot), g_1,g_2 \}
\end{equation}

Здесь $\G_1$ и $\G_2$ --- конечные циклические группы порядка $r$, $g_1$ --- образующий $\G_1$, $g_2$ --- образующий $\G_2$, и $e: \G_1 \times \G_2 \to \G_T$ --- эффективно вычислимое, невырожденное, билинейное спаривание для некоторой целевой группы $\G_T$. В реальных приложениях набор параметров обычно согласовывается заранее.

Заданные некоторые параметры Groth16, \term{протокол Groth16} (Groth16 protocol) --- это тогда четвёрка вероятностных полиномиальных алгоритмов $(\textsc{Setup},\textsc{Prove},\textsc{Vfy},\textsc{Sim})$, таких что выполняются следующие условия:
\begin{itemize}
\item (Фаза установки): $(CRS,\Tau)\leftarrow \textsc{Setup}(R)$: Алгоритм $\textsc{Setup}$ принимает R1CS $R$ на вход и вычисляет \term{\concept{общую опорную строку}} (Common Reference String, CRS) $CRS$ и \term{симуляционный люк} (simulation trapdoor) $\Tau$.
\item (Фаза доказывающего): $\pi\leftarrow \textsc{Prove}(R,CRS,I,W)$: Задано конструктивное доказательство $(I;W)$ для $R$, алгоритм $\textsc{Prove}$ принимает R1CS $R$, \concept{общую опорную строку} $CRS$ и конструктивное доказательство $(I,W)$ на вход и вычисляет zk-SNARK $\pi$.
\item (Фаза проверки): $\{\mathtt{accept},\mathtt{reject}\}\leftarrow \textsc{Vfy}(R,CRS,I,\pi)$: Алгоритм \textsc{Vfy} принимает R1CS $R$, \concept{общую опорную строку} $CRS$, экземпляр $I$ и zk-SNARK $\pi$ на вход и возвращает \texttt{reject} или \texttt{accept}.
\item (Моделирование) $\pi\leftarrow \textsc{Sim}(R,\Tau,CRS, I)$: Алгоритм \textsc{Sim} принимает R1CS $R$, \concept{общую опорную строку} $CRS$, симуляционный люк $\Tau$ и экземпляр $I$ на вход и возвращает zk-SNARK $\pi$.
\end{itemize}
Мы объясним эти алгоритмы вместе с подробными примерами в оставшейся части этого раздела.

При предположении доверенной третьей стороны или наличия соответствующего многостороннего вычисления для установки, протокол способен получить zk-SNARK из конструктивного доказательства для $R$, при условии, что порядок группы $r$ достаточно велик, причём требование особенно применимо к быть большим, чем число ограничений в ассоциированной R1CS.

\begin{example}[Задача 3-факторизации]
\label{ex:3-fac-groth-16-params} Рассмотрим задачу 3-факторизации из \ref{ex:3-factorization} и её ассоциированную алгебраическую схему \ref{ex:3-fac-zk-circuit}, а также систему ограничений ранга 1 из \ref{ex:3-factorization-r1cs}. В этом примере мы хотим согласовать набор параметров $\{r, \G_1, \G_2, e(\cdot,\cdot), g_1, g_2\}$, чтобы использовать протокол Groth16 для нашей задачи 3-факторизации.

Чтобы найти подходящие параметры, сначала заметим, что схема \ref{ex:3-fac-zk-circuit}, а также её ассоциированная R1CS $R_{3.fac\_zk}$ \ref{ex:3-factorization-r1cs} и порождённая QAP \ref{ex:3-fac-QAP}, определены над полем $\F_{13}$. Поэтому мы должны выбрать группы спаривания $\G_1$ и $\G_2$ порядка 13.

Как нам известно из \ref{BLS6}, кривая moon-math \texttt{BLS6\_6} имеет две подгруппы $\G_1[13]$ и $\G_2[13]$, которые обе имеют порядок 13. Ассоциированное спаривание Вейля $e(\cdot,\cdot)$ \ref{BLS6-weil-pairing} эффективно вычислимо, билинейно и невырождено. Поэтому мы выбираем эти группы и спаривание Вейля вместе с образующими $g_1 = (13,15)$ и $g_2=(7v^2,16v^3)$ групп $\G_1[13]$ и $\G_2[13]$ в качестве набора параметров:
$$
\mathtt{Groth16\text{-}Param}(R_{3.fac\_zk})= \{ 13, \G_1[13], \G_2[13], e(\cdot,\cdot), (13,15),(7v^2,16v^3) \}
$$
Следует отметить, что наш выбор не является единственным. Каждая пара конечных циклических групп порядка 13, имеющая эффективно вычислимое, невырожденное, билинейное спаривание, квалифицируется как набор параметров Groth16. Ситуация аналогична реальным приложениям, где SNARK с эквивалентным поведением определены над различными кривыми, используемыми в различных приложениях.
\end{example}
\begin{example}[Задача 3-факторизации в Circom и Snarkjs]
\label{ex:3-fac-groth-16-params-circom} Snark.js --- это JavaScript-библиотека, которая облегчает разработку систем, включающих доказательства с нулевым разглашением (ZKP), включая протокол Groth-16. Чтобы продемонстрировать практический пример задачи 3-факторизации, мы используем нашу реализацию на Circom (см. \ref{ex:3-fac-circom}), которая компилируется в форму, совместимую с snark.js.

На момент написания этого текста Snark.js поддерживает эллиптические кривые \curvename{alt\_bn128}, \curvename{BLS12-381} и \curvename{Goldilocks}. Для целей этого примера мы будем использовать \curvename{alt\_bn128} и её ассоциированное скалярное поле $\F_{bn128}$, как представлено в \ref{BN128}. Параметры Groth-16 для этой кривой, как официально определено для блокчейна Ethereum, можно найти в \href{https://github.com/ethereum/EIPs/blob/master/EIPS/eip-197.md}{EIP-197}. Snark.js использует эти параметры.
\end{example}
\begin{exercise}
\label{ex:baby-jubjub-circom} Реализуйте уравнение кривой Беби-ДжабДжаб (Baby-JubJub) в форме скрученной кривой Эдвардса (twisted Edwards curve) в Circom и скомпилируйте его в R1CS и ассоциированный генератор свидетелей.
\end{exercise}

\subsection{Фаза установки} Генерация zk-SNARK из конструктивных доказательств в протоколе Groth16 требует выполнения предобрабатывающей фазы. Эта фаза должна быть выполнена один раз для каждой системы ограничений ранга 1 и её ассоциированной квадратичной арифметической программы. Результатом этой фазы является \concept{общая опорная строка} (Common Reference String), которая необходима как доказывающему, так и проверяющему для генерации и проверки zk-SNARK. Дополнительно, во время этой фазы генерируется симуляционный люк, который может быть использован для моделирования доказательств.

Более точно, пусть $L$ --- язык, определённый некоторой системой ограничений ранга 1 $R$, такой что конструктивное доказательство знания для экземпляра $<I_1,\ldots,I_n>$ в $L$ состоит из свидетеля $<W_1,\ldots,W_m>$. Пусть $QAP(R) = \left\{T\in \F[x],\left\{A_j,B_j,C_j\in \F[x]\right\}_{j=0}^{n+m}\right\}$ --- квадратичная арифметическая программа (Quadratic Arithmetic Program), ассоциированная с $R$, и пусть $\{\G_1, \G_2, e(\cdot,\cdot), g_1, g_2, \F_r\}$ --- набор параметров Groth16.

Фаза установки затем выбирает 5 случайных, обратимых элементов $\alpha$, $\beta$,$\gamma$, $\delta$ и $\tau$ из скалярного поля $\F_r$ протокола и выдаёт \term{симуляционный люк} (simulation trapdoor) $\Tau$:
\begin{equation}
\label{def:groth16-trapdoor}
\Tau = (\alpha, \beta, \gamma, \delta, \tau)
\end{equation}
Кроме того, фаза установки использует эти 5 случайных элемента вместе с двумя образующими $g_1$ и $g_2$ и квадратичной арифметической программой для генерации \term{\concept{общей опорной строки}} (Common Reference String) $CRS_{QAP}= (CRS_{\mathbb{G}_1},CRS_{\mathbb{G}_2})$ языка $L$:

\begin{definition}[\deftitle{Общая \concept{опорная строка}}]
\begin{align*}
\label{def:groth16-crs}
CRS_{\mathbb{G}_{1}} &= \textstyle\left\{ \begin{array}{c}
g_1^\alpha,g_1^\beta,g_1^\delta,\left(g_1^{\tau^j},\ldots\right)_{j=0}^{deg(T)-1},
\left(g_1^{\frac{\beta\cdot A_{j}(\tau)+\alpha\cdot B_{j}(\tau)+C_{j}(\tau)}{\gamma}},\ldots\right)_{j=0}^n\\
\left(g_1^{\frac{\beta\cdot A_{j+n}(\tau)+\alpha\cdot B_{j+n}(\tau)+C_{j+n}(\tau)}{\delta}},\ldots\right)_{j=1}^m,\left(g_1^{\frac{\tau^{j}\cdot T(\tau)}{\delta}},\ldots\right)_{j=0}^{deg(T)-2}
\end{array}\right\} \\
CRS_{\mathbb{G}_{2}} &= \left\{g_2^\beta ,g_2^\gamma,g_2^\delta,\left(g_2^{\tau^j},\ldots\right) _{j=0}^{deg(T)-1}\right\}
\end{align*}
\end{definition}

Общие \concept{опорные строки} зависят от симуляционного люка и поэтому не являются единственными для задачи. Любой язык может иметь более одной \concept{общей опорной строки}. Размер \concept{общей опорной строки} линеен от размера экземпляра и размера свидетеля.

Если дан симуляционный люк $\Tau = (\alpha,\beta,\gamma,\delta, \tau)$, мы называем элемент $\tau$ \term{секретной точкой вычисления} (secret evaluation point) протокола, потому что если $\F_r$ --- скалярное поле конечных циклических групп $\G_1$ и $\G_2$, то ключевая особенность любой \concept{общей опорной строки} заключается в том, что она предоставляет данные для вычисления значения любого многочлена $P\in \F_r[x]$ степени $deg(P)<deg(T)$ в точке $\tau$ в показателе степени образующего $g_1$ или $g_2$, не зная $\tau$.

Более точно, пусть $\tau$ --- секретная точка вычисления, и пусть $P(x)=a_0\cdot x^0 + a_1\cdot x^1 + \ldots a_k\cdot x^k$ --- многочлен степени $k<deg(T)$ с коэффициентами в $\F_r$. Тогда мы можем вычислить $g_1^{P(\tau)}$, не зная фактического значения $\tau$:

\label{eq:exp_evaluation-poly}
\begin{align}
g_1^{P(\tau)} & = g_1^{a_0\cdot \tau^0 + a_1\cdot \tau^1 + \ldots a_k\cdot \tau^k} \notag\\
 & = g_1^{a_0\cdot \tau^0} \cdot g_1^{a_1\cdot \tau^1} \cdot \ldots \cdot g_1^{a_k\cdot \tau^k}\notag\\
 & = \Big(g_1^{\tau^0}\Big)^{a_0} \cdot \Big(g_1^{\tau^1}\Big)^{a_1} \cdot \ldots \cdot \Big(g_1^{\tau^k}\Big)^{a_k}
\end{align}
В этом выражении все групповые точки $g_1^{\tau^j}$ являются частью \concept{общей опорной строки}, следовательно, они могут быть использованы для вычисления результата. То же самое верно для вычисления $g_2^{P(\tau)}$, поскольку $\G_2$ часть \concept{общей опорной строки} содержит точки $g_2^{\tau^j}$.

В практических приложениях элементы $g_{1/2}^{\tau^0}$, $g_{1/2}^{\tau^1}$, $\ldots$, $g_{1/2}^{\tau^k}$ обычно называются \term{степенями тау} (powers of tau), термин, часто используемый в вычислениях доверенной установки. Кроме того, симуляционный люк часто называют \term{токсичными отходами} (toxic waste) фазы установки. Как будет показано в разделе \ref{sec:proof_simulation}, симуляционный люк может быть использован для генерации мошеннических доказательств, которые являются проверяемыми zk-SNARK, которые могут быть построены без знания какого-либо свидетеля. Общая опорная строка также известна как \term{пара ключей доказывающего и проверяющего} (prover and verifier key pair).

Чтобы обеспечить безопасность протокола, установка должна быть выполнена таким образом, который гарантирует безопасное уничтожение симуляционного люка. Простейший метод достижения этого --- использование так называемой \term{доверенной третьей стороны} (trusted third party), где доверие возлагается на сторону должным образом сгенерировать \concept{общую опорную строку} и безопасно уничтожить токсичные отходы после этого.

Однако найти доверенную третью сторону может быть сложно, поэтому в практических приложениях были разработаны альтернативные методы. Они используют многостороннее вычисление на фазе установки, которое может быть публично проверено на правильное выполнение, и симуляционный люк не может быть восстановлен, если хотя бы один участник уничтожит свой вклад. Каждый участник держит лишь часть люка, делая его восстановимым только если все участники сотрудничают и делятся своими частями.

\begin{example}[Задача 3-факторизации]
\label{ex:3-fac-groth-16-crs}
Чтобы увидеть, как может быть вычислена фаза установки Groth16 zk-SNARK, рассмотрим задачу 3-факторизации из \ref{ex:3-factorization} и параметры Groth16 из \examplename{} \ref{ex:3-fac-groth-16-params}. Как мы видели в \ref{ex:3-fac-QAP}, ассоциированная квадратичная арифметическая программа задаётся следующим множеством:
\begin{multline*}
QAP(R_{3.fac\_zk}) =\{x^{2}+x+9,\\
 \{0,0,6x+10,0,0,7x+4\},\{0,0,0,6x+10,7x+4,0\},\{0,7x+4,0,0,0,6x+10\}\}
\end{multline*}
Чтобы преобразовать эту QAP в \concept{общую опорную строку}, мы выбираем элементы поля $\alpha=6$, $\beta=5$, $\gamma=4$, $\delta=3$, $\tau=2$ из $\mathbb{F}_{13}$.
В реальных приложениях важно выбирать эти значения случайным образом из скалярного поля, но в нашем подходе мы выбираем эти неслучайные значения, чтобы сделать их более запоминающимися, что помогает при вычислениях на бумаге. Наш симуляционный люк затем задаётся следующим образом:
$$
\Tau = (6,5,4,3,2)
$$

Мы храним это в секрете, чтобы моделировать доказательства позже, но мы осторожно скрываем $\Tau$ от всех, кто не читал эту книгу. Затем мы инстанцируем \concept{общую опорную строку} \ref{def:groth16-crs} из этих значений. Поскольку наши группы являются подгруппами эллиптической кривой \texttt{BLS6\_6}, мы используем обозначение скалярного произведения вместо возведения в степень.

Чтобы вычислить $\G_1$ часть \concept{общей опорной строки}, мы используем логарифмический порядок группы $\G_1$ \ref{BLS6-G1-log}, образующий $g_1=(13,15)$, а также значения из симуляционного люка. Поскольку $deg(T)=2$, мы получаем следующее:
\begin{align*}
[\alpha]g_1 & = [6](13,15) = (27,34) \\
[\beta]g_1 & = [5](13,15) = (26,34) \\
[\delta]g_1 & = [3](13,15) = (38,15)
\end{align*}
Чтобы вычислить остальную $\G_1$ часть \concept{общей опорной строки}, мы раскрываем индексированные кортежи и подставляем секретные случайные элементы из симуляционного люка. Мы получаем следующее:
\begin{align*}
\Big( [\tau^{j}]g_1,\ldots\Big) _{j=0}^{1} = 
 & \Big( [2^0](13,15), [2^1](13,15)\Big)  \\
 = & \Big((13,15),(33,34)\Big)\\
\Big([\frac{\beta A_{j}(\tau)+ \alpha B_{j}(\tau) + C_{j}(\tau)}{\gamma}]g_1,\ldots\Big)_{j=0}^1 =
 & \Big([\frac{5 A_{0}(2)+6 B_{0}(2)+C_{0}(2)}{4}](13,15),\\
 &\phantom{\Big(} [\frac{5 A_{1}(2)+6 B_{1}(2)+C_{1}(2)}{4}](13,15)\Big)\\
\Big([\frac{\beta A_{j+n}(\tau)+ \alpha B_{j+n}(\tau) + C_{j+n}(\tau)}{\delta}]g_1,\ldots\Big)_{j=1}^4 = 
&  \Big( [\frac{5 A_{2}(2)+ 6 B_{2}(2) + C_{2}(2)}{3}](13,15),\\ 
& \phantom{\Big(} [\frac{5 A_{3}(2)+ 6 B_{3}(2) + C_{3}(2)}{3}](13,15),\\
& \phantom{\Big(} [\frac{5 A_{4}(2)+ 6 B_{4}(2) + C_{4}(2)}{3}](13,15),\\ 
& \phantom{\Big(} [\frac{5 A_{5}(2)+ 6 B_{5}(2) + C_{5}(2)}{3}](13,15)\Big)\\
\Big([\frac{\tau^j\cdot T(\tau)}{\delta})]g_1\Big)_{j=0}^0 = & \Big([\frac{2^0\cdot T(2)}{3}](13,15)\Big) 
\end{align*}
Чтобы вычислить точки кривой в правой части этих выражений, нам нужны многочлены из ассоциированной квадратичной арифметической программы и вычислить их в секретной точке $\tau=2$. Поскольку $4^{-1}=10$ и $3^{-1}=9$ в $\F_{13}$, мы получаем следующее:
\begin{align*}
[\frac{5 A_{0}(2)+6 B_{0}(2)+C_{0}(2)}{4}](13,15) = 
 & [(5 \cdot 0 +6\cdot 0 + 0)\cdot 10](13,15) = [0](13,14)\\
 & \mathcal{O} \\
[\frac{5 A_{1}(2)+6 B_{1}(2)+C_{1}(2)}{4}](13,15) = 
 & [(5\cdot 0 +6\cdot 0 + (7\cdot 2 + 4))\cdot 10](13,15) = [11](13,15) = \\
 & (33,9) \\
[\frac{5 A_{2}(2)+ 6 B_{2}(2) + C_{2}(2)}{3}](13,15) =
 & [(5\cdot (6\cdot 2 +10) +6\cdot 0 +0 )\cdot 9](13,15) = [2](13,15) = \\
 & (33,34) \\
[\frac{5 A_{3}(2)+ 6 B_{3}(2) + C_{3}(2)}{3}](13,15) =
 & [(5\cdot 0 +6\cdot (6\cdot 2 + 10) + 0 )\cdot 9](13,15) = [5](13,15) = \\
 & (26,34) \\
[\frac{5 A_{4}(2)+ 6 B_{4}(2) + C_{4}(2)}{3}](13,15) = 
 & [(5\cdot 0+6\cdot(7\cdot 2 +4)+0)\cdot 9](13,15) =[10](13,15) = \\
 & (38,28) \\
[\frac{5 A_{5}(2)+ 6 B_{5}(2) + C_{5}(2)}{3}](13,15) =
 & [(5\cdot (7\cdot 2 + 4) +6\cdot 0 + 6\cdot 2 + 10 )\cdot 9](13,15) = [7](13,15) = \\
 & (27,9)\\
[\frac{2^0\cdot T(2)}{3}](13,15) =
 & [1\cdot (2^2+2+9)\cdot 9](13,15)= [5](13,15) = \\
 & (26,34)
\end{align*}
Собрав все эти значения вместе, мы видим, что $\mathbb{G}_1$ часть \concept{общей опорной строки} (Common Reference String) задаётся следующим множеством из $12$ точек из $13$-кручённой группы $\G_1$ кривой \texttt{BLS6\_6}:
\begin{equation}
\label{eq:3-fac-groth-16-crs}
CRS_{\mathbb{G}_{1}}=\left\{ \begin{array}{c}
(27,34),(26,34),(38,15),\Big((13,15),(33,34)\Big),
\Big(\mathcal{O}, (33,9)\Big)\\
\Big((33,34),(26,34),(38,28),(27,9)\Big),
\Big((26,34)\Big)
\end{array}\right\}
\end{equation}

Чтобы вычислить $\G_2$ часть \concept{общей опорной строки}, мы используем логарифмический порядок группы $\G_2$ \ref{BLS6-G2-log}, образующий $g_2=(7v^2,16v^3)$, а также значения из симуляционного люка. Поскольку $deg(T)=2$, мы получаем следующее:
\begin{align*}
[\beta]g_2 & = [5](7v^2,16v^3) = (16v^2,28v^3) \\
[\gamma]g_2 & = [4](7v^2,16v^3) = (37v^2,27v^3) \\
[\delta]g_2 & = [3](7v^2,16v^3) = (42v^2,16v^3)
\end{align*}
Чтобы вычислить остальную $\G_2$ часть \concept{общей опорной строки}, мы раскрываем индексированный кортеж и подставляем секретные случайные элементы из симуляционного люка. Мы получаем следующее:
\begin{align*}
\Big( [\tau^{j}]g_2,\ldots\Big) _{j=0}^{1} = 
 & \Big( [2^0](7v^2,16v^3), [2^1](7v^2,16v^3)\Big)  \\
 = & \Big((7v^2,16v^3),(10v^2,28v^3)\Big)
\end{align*}
Собрав все эти значения вместе, мы видим, что $\mathbb{G}_2$ часть \concept{общей опорной строки} задаётся следующим множеством из $5$ точек из $13$-кручённой группы $\G_2$ кривой \texttt{BLS6\_6}:
$$
CRS_{\mathbb{G}_{2}}=\left\{(16v^2,28v^3) ,(37v^2,27v^3),(42v^2,16v^3),\Big(7v^2,16v^3), (10v^2,28v^3)\Big)\right\} 
$$
Заданы симуляционный люк $\Tau$ и квадратичная арифметическая программа \ref{ex:3-fac-QAP}, ассоциированная \concept{общая опорная строка} задачи 3-факторизации выглядит следующим образом:
\begin{align*}
CRS_{\mathbb{G}_{1}} &=\left\{ \begin{array}{c}
(27,34),(26,34),(38,15),\Big((13,15),(33,34)\Big),
\Big(\mathcal{O}, (33,9)\Big)\\
\Big((33,34),(26,34),(38,28),(27,9)\Big),
\Big((26,34)\Big)
\end{array}\right\}\\
CRS_{\mathbb{G}_{2}} &=\left\{(16v^2,28v^3) ,(37v^2,27v^3),(42v^2,16v^3),\Big(7v^2,16v^3), (10v^2,28v^3)\Big)\right\}
\end{align*}
Мы затем публикуем эти данные для всех, кто хочет участвовать в генерации zk-SNARK или его проверке для задачи 3-факторизации.

Чтобы понять, как эта \concept{общая опорная строка} может быть использована для вычисления многочленов в секретной точке вычисления в показателе степени образующего, предположим, что мы удалили симуляционный люк. В этом случае, при предположении, что проблема дискретного логарифма является трудной в наших группах, у нас нет способа узнать секретную точку вычисления, следовательно, мы не можем вычислять многочлены в этой точке. Однако мы можем вычислять многочлены меньшей степени, чем степень целевого многочлена, в показателе степени обоих образующих в этой точке.

Чтобы увидеть это, рассмотрим, например, многочлены $A_2(x)= 6x +10$ и $A_5(x)=7x+4$ из QAP этой задачи. Чтобы вычислить эти многочлены в показателе степени $g_1$ и $g_2$ в секретной точке $\tau$, не зная значения $\tau$ (которое в нашем случае $2$), мы можем использовать \concept{общую опорную строку} и уравнение \ref{eq:exp_evaluation-poly}. Используя обозначение скалярного произведения вместо возведения в степень, мы получаем следующее:
\begin{align}
[A_2(\tau)]g_1 & = [6\cdot \tau^1 + 10\cdot \tau^0] g_1 \nonumber \\
     & = [6](33,34) + [10](13,15) & \text{\# } [\tau^0]g_1 = (13,15), [\tau^1]g_1 = (33,34) \nonumber \\
     & = [6\cdot 2](13,15) + [10](13,15) = [9](13,15) & \text{\# логарифмический порядок на } \G_1 \nonumber \\
     & = (35,15) \label{eq:3-fac-A2-g1} \\
[A_5(\tau)]g_1 & = [7\cdot \tau^1 + 4\cdot \tau^0] g_1 \nonumber \\
     & = [7](33,34) + [4](13,15) \nonumber \\
     & = [7\cdot 2](13,15) + [4](13,15) = [5](13,15) \nonumber \\
     & = (26,34) \label{eq:3-fac-A5-g1}
\end{align}

Действительно, мы способны вычислять многочлены в показателе степени в секретной точке вычисления, потому что эта точка зашифрована в точке кривой $(33,34)$, и её секретность защищена предположением о дискретном логарифме. Конечно, в нашем вычислении мы восстановили секретную точку $\tau=2$, но это было возможно только потому, что мы знаем логарифмический порядок наших групп относительно образующих. Такой порядок является неосуществимым в криптографически защищённых кривых. Мы можем выполнить то же вычисление на $\G_2$ и получить следующее:
\begin{align}
[A_2(\tau)]g_2 & = [6\cdot \tau^1 + 10\cdot \tau^0] g_2 \nonumber \\
     & = [6](10v^2,28v^3) + [10](7v^2,16v^3) \nonumber \\
     & = [6\cdot 2](7v^2,16v^3) + [10](7v^2,16v^3) = [9](7v^2,16v^3) \nonumber \\
     & = (37v^2,16v^3) \label{eq:3-fac-A2-g2} \\
[A_5(\tau)]g_2 & = [7\cdot \tau^1 + 4\cdot \tau^0] g_1 \nonumber \\
     & = [7](10v^2,28v^3) + [4](7v^2,16v^3) \nonumber \\
     & = [7\cdot 2](7v^2,16v^3) + [4](7v^2,16v^3) = [5](7v^2,16v^3) \nonumber \\
     & = (16v^2,28v^3)  \label{eq:3-fac-A5-g2}
\end{align}

Помимо целевого многочлена $T$, все остальные многочлены квадратичной арифметической программы могут быть вычислены в показателе степени таким образом.
\end{example}

\begin{example}[Задача 3-факторизации в Circom и Snark.js]
\label{ex:3-fac-groth-16-setup-circom} Реализацию фазы установки Groth16 zk-SNARK в реальных приложениях можно наблюдать посредством изучения нашей реализации задачи 3-факторизации на Circom из \ref{ex:3-fac-circom} и ассоциированного набора параметров из Snark.js, как описано в примере \ref{ex:3-fac-groth-16-params-circom}.

В соответствии с методологией, описанной в \cite{bowe-17}, генерация общей опорной строки (Common Reference String) в Snark.js состоит из двух частей. Первая часть зависит от верхней границы числа ограничений в схеме, в то время как вторая часть зависит от самой схемы. Это разделение повышает гибкость процедуры доверенной установки, поскольку протоколы с универсальными общими опорными строками, такие как PLONK, требуют выполнения только первой фазы, в то время как протокол Groth16 требует выполнения обеих фаз.

Первая фаза, обычно называемая \term{степенями тау} (powers of tau), включает вычисление последовательных степеней, $g^\tau$, $g^{\tau^2}$, $g^{\tau^3}$, $\ldots$, $g^{\tau^k}$, случайного элемента поля $\tau$ в показателях степени согласованных образующих $\G_1$ и $\G_2$. Случайный элемент $\tau$ может быть либо предоставлен доверенной третьей стороной, либо сгенерирован через многостороннее вычисление.

Предполагая, что верхняя граница числа ограничений в схеме 3-fac Circom задана $2^4$, первая часть инициализируется следующим образом:
\\
\\
\texttt{:\$ snarkjs powersoftau new bn128 4 pot4\_0000.ptau -v}
\\
\\
Команда \texttt{new} используется для запуска новой фазы, и первый параметр после \texttt{new} относится к типу кривой и, следовательно, задаёт набор параметров Groth-16, как определено в \ref{ex:3-fac-groth-16-params-circom}. Следующий параметр \texttt{a}, в данном случае 4, определяет верхнюю границу $2^a$ числа ограничений, которые установка может принять.

После этой инициализации несколько сторон могут внести случайность в общую опорную строку:
\\
\\
\texttt{:\$ snarkjs powersoftau contribute pot4\_0000.ptau pot4\_0001.ptau \-\-name="1st\_cont" -v}
\\
\\
Этот шаг может повторяться, и вклад случайности каждой стороной в общую опорную строку облегчается созданием нового файла вклада, и пользователю предлагается ввести дополнительную энтропию. Было показано, что при условии, что хотя бы один участник забудет свою случайность, восстановление симуляционного люка, определённого в ссылке \ref{def:groth16-trapdoor}, невозможно. Проверка корректности любого файла вклада выполняется следующим образом:
\\
\\
\texttt{:\$ snarkjs powersoftau verify pot4\_0001.ptau}
\\
\\
Завершение первой фазы требует включения некоторой публичной и непредсказуемой случайности в генерацию степеней тау, которая должна быть непредсказуемой до тех пор, пока не сделан вклад последнего участника. Это обычно достигается использованием хэша последнего блока в публично доступном блокчейне или аналогичного механизма.
\\
\\
\texttt{:\$ snarkjs powersoftau beacon pot4\_0001.ptau pot4\_beacon.ptau \ 0102030405060708090a0b0c00 10 -n="Final Beacon"}
\\
\\
После включения вклада случайности от участвующих сторон и случайного маяка следующие две команды служат для завершения и проверки первой фазы:
\\
\\
\texttt{:\$ snarkjs powersoftau prepare phase2 pot4\_beacon.ptau \ pot4\_final.ptau -v}\\
\texttt{:\$ snarkjs powersoftau verify pot4\_final.ptau}
\\
\\
Очевидно, что эта фаза зависит только от верхней границы числа ограничений и не зависит от какой-либо конкретной схемы. Следовательно, эта церемония степеней тау может быть использована для генерации общей опорной строки для любой схемы с числом ограничений меньше $2^4$.

Как указано в ссылке \ref{ex:3-fac-circom}, схема three\_fac в Circom состоит из двух ограничений. Следовательно, случайность, сгенерированная из этой церемонии степеней тау, может быть использована для инициации второй фазы процесса генерации общей опорной строки:
\\
\\
\texttt{:\$ snarkjs groth16 setup three\_fac.r1cs pot4\_final.ptau \ three\_fac0000.zkey}
\\
\\
Вторая фаза начинается с вычисления общей опорной строки Groth16, как указано в ссылке \ref{def:groth16-crs}. Эта фаза зависит от R1CS задачи, и результирующее вычисление сохраняется в \fname{three\_fac0000.zkey}.

В отличие от первой фазы, где случайность использовалась для генерации параметра '$\tau$' CRS, несколько сторон теперь могут внести случайность в другие параметры, а именно $\alpha$, $\beta$, $\gamma$ и $\delta$, CRS Groth16.
\\
\\
\texttt{:\$ snarkjs zkey contribute three\_fac0000.zkey \ three\_fac0001.zkey --name="1st Contributor Name" -v}
\\
\\
Это создаёт файл с новым вкладом во вторую фазу и предлагает пользователю предоставить дополнительную случайность для повышения энтропии. Этот шаг может повторяться много раз для разных пользователей. Корректность любого файла вклада может быть проверена следующим образом:
\\
\\
\texttt{:\$ snarkjs zkey verify three\_fac.r1cs pot4\_final.ptau \ three\_fac0001.zkey}
\\
\\
Вторая фаза требует интеграции непредсказуемой, публичной случайности, которая не известна до вклада последнего участника. Это может быть достигнуто посредством использования такого метода, как хэш последнего блока в публичном блокчейне:
\\
\\
\texttt{:\$ snarkjs zkey beacon three\_fac0001.zkey three\_fac\_final.zkey \
010203040506070809 10 -n="Final Beacon phase2"}
\\
\\
После того как вклад случайности от всех сторон и случайный маяк были включены, проверка второй фазы может быть выполнена с использованием следующих двух команд. Этот процесс также экспортирует ключ проверки, который является жизненно важным компонентом общей опорной строки и будет сохранён в JSON-файле для использования проверяющим:
\\
\\
\texttt{:\$ snarkjs zkey verify three\_fac.r1cs pot4\_final.ptau \ three\_fac\_final.zkey}\\
\texttt{:\$ snarkjs zkey export verificationkey three\_fac\_final.zkey \ verification\_key.json}
\\
\\
Доверенная установка Circom, таким образом, позволяет создать общую опорную строку \fname{three\_fac\_final.zkey}, которая содержит как ключ доказывающего, так и ключ проверяющего. Ключ проверяющего также может быть экспортирован как JSON-файл \fname{verification\_key.json} с целью публикации его в публичном блокчейне для реализации в виде смарт-контракта.
\end{example}
\begin{exercise} 
\label{ex:baby-jubjub-circom-setupt}
Рассмотрите упражнение \ref{ex:baby-jubjub-circom} и выполните 3-стороннюю фазу доверенной установки для схемы baby-jubjub.
\end{exercise}

\subsection{Фаза доказывающего} Чтобы понять, как общая опорная строка может быть использована в фазе доказывающего для генерации zk-SNARK, давайте повторно рассмотрим пример задачи 3-факторизации и предположим, что в дополнение к общей опорной строке доказывающий выбрал некоторое значение $I_1=11$. Тогда доказывающий должен вычислить допустимое назначение $W$ для экземпляра $<3\_fac,3\_fac.r1cs, I_1>$ из $\mathcal{R}_{3\_fac.r1cs}$. В данном случае $W=(11,3,2,6,1)$, что является допустимым назначением, поскольку $11\cdot 3\cdot 2=66$ и остальные строки таблицы \ref{ex:3-fac-qap-witness} выполняются при $(I_1;W) = (11;3,2,6,1)$. Доказывающий использует секретное знание $W$ в качестве свидетеля и выполняет фазу доказывающего алгоритма Groth16 zk-SNARK, описанную в ссылке \ref{groth16}. 

Чтобы вычислить открытую часть $S_1$ из определения \ref{def:groth16-snark}, доказывающий должен сначала вычислить элементы $\F_{13}$, заданные следующими полиномами и их коэффициентами:
\begin{align*}
P_1(x) & = l_1(x)\cdot 11 + l_3(x)\cdot 3 + l_5(x)\cdot 2 + l_6(x)\cdot 6 + l_7(x)\cdot 1 \\
       & = 10\cdot 11 + 2\cdot 3 + 4\cdot 2 + 8\cdot 6 + 1\cdot 1 = 11\\
P_2(x) & = l_1(x)\cdot 11 + l_3(x)\cdot 3 + l_5(x)\cdot 2 + l_6(x)\cdot 6 + l_7(x)\cdot 1 \\
       & = 10\cdot 11 + 5\cdot 3 + 10\cdot 2 + 5\cdot 6 + 1\cdot 1 = 11\\
P_3(x) & = l_1(x)\cdot 11 + l_3(x)\cdot 3 + l_5(x)\cdot 2 + l_6(x)\cdot 6 + l_7(x)\cdot 1 \\
       & = 10\cdot 11 + 1\cdot 3 + 2\cdot 2 + 4\cdot 6 + 1\cdot 1 = 9\\
P_4(x) & = l_1(x)\cdot 11 + l_3(x)\cdot 3 + l_5(x)\cdot 2 + l_6(x)\cdot 6 + l_7(x)\cdot 1 \\
       & = 10\cdot 11 + 1\cdot 3 + 4\cdot 2 + 8\cdot 6 + 1\cdot 1 = 9\\
P_5(x) & = l_1(x)\cdot 11 + l_3(x)\cdot 3 + l_5(x)\cdot 2 + l_6(x)\cdot 6 + l_7(x)\cdot 1 \\
       & = 10\cdot 11 + 1\cdot 3 + 1\cdot 2 + 1\cdot 6 + 1\cdot 1 = 10\\
P_6(x) & = l_1(x)\cdot 11 + l_3(x)\cdot 3 + l_5(x)\cdot 2 + l_6(x)\cdot 6 + l_7(x)\cdot 1 \\
       & = 10\cdot 11 + 1\cdot 3 + 1\cdot 2 + 1\cdot 6 + 1\cdot 1 = 10
\end{align*}
и
\begin{align*}
R_1(x) & = r_1(x)\cdot 11 + r_3(x)\cdot 3 + r_5(x)\cdot 2 + r_6(x)\cdot 6 + r_7(x)\cdot 1 \\
       & = 10\cdot 11 + 1\cdot 3 + 2\cdot 2 + 4\cdot 6 + 1\cdot 1 = 9\\
R_2(x) & = r_1(x)\cdot 11 + r_3(x)\cdot 3 + r_5(x)\cdot 2 + r_6(x)\cdot 6 + r_7(x)\cdot 1 \\
       & = 10\cdot 11 + 4\cdot 3 + 1\cdot 2 + 2\cdot 6 + 1\cdot 1 = 9\\
R_3(x) & = r_1(x)\cdot 11 + r_3(x)\cdot 3 + r_5(x)\cdot 2 + r_6(x)\cdot 6 + r_7(x)\cdot 1 \\
       & = 10\cdot 11 + 1\cdot 3 + 1\cdot 2 + 1\cdot 6 + 1\cdot 1 = 10\\
R_4(x) & = r_1(x)\cdot 11 + r_3(x)\cdot 3 + r_5(x)\cdot 2 + r_6(x)\cdot 6 + r_7(x)\cdot 1 \\
       & = 10\cdot 11 + 1\cdot 3 + 1\cdot 2 + 1\cdot 6 + 1\cdot 1 = 10\\
R_5(x) & = r_1(x)\cdot 11 + r_3(x)\cdot 3 + r_5(x)\cdot 2 + r_6(x)\cdot 6 + r_7(x)\cdot 1 \\
       & = 10\cdot 11 + 1\cdot 3 + 1\cdot 2 + 1\cdot 6 + 1\cdot 1 = 10\\
R_6(x) & = r_1(x)\cdot 11 + r_3(x)\cdot 3 + r_5(x)\cdot 2 + r_6(x)\cdot 6 + r_7(x)\cdot 1 \\
       & = 10\cdot 11 + 1\cdot 3 + 1\cdot 2 + 1\cdot 6 + 1\cdot 1 = 10
\end{align*}
и
\begin{align*}
O_1(x) & = o_1(x)\cdot 11 + o_3(x)\cdot 3 + o_5(x)\cdot 2 + o_6(x)\cdot 6 + o_7(x)\cdot 1 \\
       & = 10\cdot 11 + 1\cdot 3 + 1\cdot 2 + 1\cdot 6 + 1\cdot 1 = 10\\
O_2(x) & = o_1(x)\cdot 11 + o_3(x)\cdot 3 + o_5(x)\cdot 2 + o_6(x)\cdot 6 + o_7(x)\cdot 1 \\
       & = 10\cdot 11 + 1\cdot 3 + 1\cdot 2 + 1\cdot 6 + 1\cdot 1 = 10\\
O_3(x) & = o_1(x)\cdot 11 + o_3(x)\cdot 3 + o_5(x)\cdot 2 + o_6(x)\cdot 6 + o_7(x)\cdot 1 \\
       & = 10\cdot 11 + 4\cdot 3 + 8\cdot 2 + 1\cdot 6 + 1\cdot 1 = 9\\
O_4(x) & = o_1(x)\cdot 11 + o_3(x)\cdot 3 + o_5(x)\cdot 2 + o_6(x)\cdot 6 + o_7(x)\cdot 1 \\
       & = 10\cdot 11 + 1\cdot 3 + 2\cdot 2 + 4\cdot 6 + 1\cdot 1 = 11\\
O_5(x) & = o_1(x)\cdot 11 + o_3(x)\cdot 3 + o_5(x)\cdot 2 + o_6(x)\cdot 6 + o_7(x)\cdot 1 \\
       & = 10\cdot 11 + 1\cdot 3 + 4\cdot 2 + 8\cdot 6 + 1\cdot 1 = 9\\
O_6(x) & = o_1(x)\cdot 11 + o_3(x)\cdot 3 + o_5(x)\cdot 2 + o_6(x)\cdot 6 + o_7(x)\cdot 1 \\
       & = 10\cdot 11 + 1\cdot 3 + 1\cdot 2 + 1\cdot 6 + 1\cdot 1 = 10
\end{align*}

Теперь, чтобы вычислить элементы $\F_{13}$, заданные этими полиномами, доказывающий должен оценить каждый из этих полиномов на всех шести элементах $T=\{1,2,3,4,5,6\}$ из определения \ref{def:3-fac-zk-circuit}. Это означает, что доказывающий должен вычислить коэффициенты $(a_0,a_1,\ldots,a_5)$ полинома $P(x)=a_0+a_1x+a_2x^2+a_3x^3+a_4x^4+a_5x^5$, который удовлетворяет $P(1)=11$, $P(2)=11$, $P(3)=9$, $P(4)=9$, $P(5)=10$ и $P(6)=10$. Это можно сделать, например, путём решения системы линейных уравнений или с помощью интерполяционной формулы Лагранжа. В любом случае мы получаем следующие полиномы:
\begin{align*}
P(x) & = 6x^5 + 8x^4 + 8x^3 + 10x^2 + 7x + 11\\
R(x) & = 8x^5 + 2x^4 + 7x^3 + 4x^2 + 9x + 11\\
O(x) & = 6x^5 + 8x^4 + 8x^3 + 10x^2 + 7x + 11
\end{align*}

Чтобы вычислить первую составляющую $g^\alpha$ из ссылки \ref{def:groth16-snark}, доказывающий должен сгенерировать два случайных поля элемента $r$ и $s$ из $\F_{13}$. Предположим, что он выбирает $r=2$ и $s=3$. Тогда доказывающий может вычислить первую составляющую $S_1$ следующим образом:
\begin{align*}
[\mathbb{A}]_1 & = g^\alpha\cdot g^{P(\tau)} \cdot g^{r\cdot \delta}\\
               & = 38^{2}\cdot 38^{P(4)}\cdot 38^{2\cdot 7}\\
               & = 38^{2}\cdot 38^{6\cdot 4^5 + 8\cdot 4^4 + 8\cdot 4^3 + 10\cdot 4^2 + 7\cdot 4 + 11}\cdot 38^{14}\\
               & = 38^{2}\cdot 38^{6\cdot 10 + 8\cdot 9 + 8\cdot 12 + 10\cdot 3 + 7\cdot 4 + 11}\cdot 38^{14}\\
               & = 38^{2}\cdot 38^{8}\cdot 38^{1} = 38^{11}\\
               & = 13
\end{align*}

Точно так же вычисляется вторая составляющая $S_1$:
\begin{align*}
[\mathbb{B}]_2 & = h^\beta \cdot h^{R(\tau)} \cdot h^{s\cdot \delta}\\
               & = 27^{5}\cdot 27^{R(4)}\cdot 27^{3\cdot 7}\\
               & = 27^{5}\cdot 27^{8\cdot 4^5 + 2\cdot 4^4 + 7\cdot 4^3 + 4\cdot 4^2 + 9\cdot 4 + 11}\cdot 27^{21}\\
               & = 27^{5}\cdot 27^{8\cdot 10 + 2\cdot 9 + 7\cdot 12 + 4\cdot 3 + 9\cdot 4 + 11}\cdot 27^{8}\\
               & = 27^{5}\cdot 27^{11}\cdot 27^{8} = 27^{11}\\
               & = 16
\end{align*}

Наконец, чтобы вычислить третью составляющую $S_1$, доказывающий должен сначала вычислить так называемый \term{факторизующий полином} $T(x)=(x-1)(x-2)(x-3)(x-4)(x-5)(x-6)$, а затем полином $H(x)$, такой что $P(x)\cdot R(x) - O(x) = T(x)\cdot H(x)$. В нашем случае мы получаем:
\begin{align*}
T(x) & = (x-1)(x-2)(x-3)(x-4)(x-5)(x-6)\\
     & = x^6 + 7x^5 + 3x^4 + 11x^3 + 10x^2 + 6x + 5
\end{align*}
и, используя алгоритм деления многочленов, мы находим $H(x)=5x^4 + 7x^3 + 9x^2 + 2x + 1$. После этого доказывающий может вычислить третью составляющую $S_1$ следующим образом:
\begin{align*}
[\mathbb{C}]_1 & = g^{H(\tau)}\cdot (g^{\beta\cdot P(\tau)+\alpha\cdot R(\tau)+O(\tau)})^1 \cdot (g^r)^\beta \cdot (g^s)^\alpha \cdot g^{r\cdot s\cdot \delta}\\
               & = 38^{H(4)}\cdot 38^{5\cdot P(4)+2\cdot R(4)+O(4)} \cdot 38^{2\cdot 5}\cdot 38^{3\cdot 2}\cdot 38^{2\cdot 3\cdot 7}\\
               & = 38^{5\cdot 10 + 7\cdot 9 + 9\cdot 3 + 2\cdot 10 + 1} \cdot 38^{5\cdot 8 + 2\cdot 8 + 8} \cdot 38^{10} \cdot 38^{6} \cdot 38^{16}\\
               & = 38^{3}\cdot 38^{5}\cdot 38^{10}\cdot 38^{6}\cdot 38^{16} = 38^{7}\\
               & = 27
\end{align*}
Таким образом, вычисляется следующий zk-SNARK:
\begin{equation}
S_{<3\_fac,3\_fac.r1cs,11>} = (13; 16; 27)
\end{equation}

\begin{example}[Схема 3-факторизации Groth16, Фаза доказывающего]
\label{ex:3-fac-groth-16-prover-circom}
Реализацию фазы доказывающего Groth16 zk-SNARK в реальных приложениях можно изучить, рассмотрев нашу реализацию задачи 3-факторизации в Circom из ссылки \ref{ex:3-fac-circom}, ассоциированные R1CS из ссылки \ref{ex:3-fac-r1cs-circom}, параметры из примера \ref{ex:3-fac-groth-16-params-circom} и общую опорную строку из примера \ref{ex:3-fac-groth-16-setup-circom}.

Предположим, что доказывающий выбрал $I_1=11$ в качестве своего секретного входа, тогда доказывающий вычисляет следующий JSON-файл свидетеля, содержащий экземпляр и назначение свидетеля $W$ в сигнальной переменной \texttt{three\_fac.in}:
\\
\\
\texttt{\{ "in": 11 \}}
\\
\\
Этот файл используется для генерации полного свидетеля, ассоциированного со схемой 3-факторизации Circom, используя следующую команду:
\\
\\
\texttt{:\$ snarkjs calculatewitness --wasm three\_fac\_js/three\_fac.wasm --input input.json --witness witness.wtns}
\\
\\
В результате генерируется файл свидетеля \fname{witness.wtns}, который включает входы $I_1$ и назначение свидетеля $W=(I_2,I_3,A,B,\mathit{out})$.

Для генерации доказательства доказывающий использует команду \texttt{groth16 prove}, используя общую опорную строку \fname{three\_fac\_final.zkey} и полный свидетель:
\\
\\
\texttt{:\$ snarkjs groth16 prove three\_fac\_final.zkey witness.wtns proof.json public.json}
\\
\\
Параметр \texttt{three\_fac\_final.zkey} содержит общую опорную строку, созданную в фазе установки. Файл \fname{proof.json} содержит zk-SNARK $(\mathbb{A},\mathbb{B},\mathbb{C})$, а \fname{public.json} содержит открытую часть экземпляра.
\end{example}
\begin{exercise}
Рассмотрите упражнение \ref{ex:baby-jubjub-circom-setupt}. Выберите некоторый секретный вход, затем вычислите соответствующее назначение свидетеля и сгенерируйте zk-SNARK.
\end{exercise}
\subsection{Фаза верификации}
Чтобы понять, как zk-SNARK может быть проверен в фазе проверки, давайте повторно рассмотрим пример задачи 3-факторизации и предположим, что некоторый проверяющий получил zk-SNARK, данный следующим образом:
\begin{equation}
S_{<3\_fac,3\_fac.r1cs,11>} = (13; 16; 27)
\end{equation}

Чтобы проверить этот zk-SNARK, проверяющий должен знать экземпляр $I_1=11$ и общую опорную строку из примера \ref{ex:groth16-crs}. Затем он вычисляет следующее:
\begin{align*}
e(\mathbb{A},\mathbb{B}) & = e(13,16) \\
                           & = e(38^{11},27^{11})\\
                           & = e(38,27)^{11\cdot 11}\\
                           & = e(38,27)^{4}\\
                           & = 9^4 = 9
\end{align*}

Чтобы вычислить правую часть парного уравнения, проверяющий должен сначала вычислить элементы $\F_{13}$, заданные полиномами $C_j(x)$ из определения \ref{def:3-fac-qap} при $j\in \{1,2,3,4\}$. Получаем:
\begin{align*}
C_1(x) & = c_1(x)\cdot 11 + c_3(x)\cdot 3 + c_5(x)\cdot 2 + c_6(x)\cdot 6 + c_7(x)\cdot 1 \\
       & = 10\cdot 11 + 1\cdot 3 + 1\cdot 2 + 1\cdot 6 + 1\cdot 1 = 10\\
C_2(x) & = c_1(x)\cdot 11 + c_3(x)\cdot 3 + c_5(x)\cdot 2 + c_6(x)\cdot 6 + c_7(x)\cdot 1 \\
       & = 10\cdot 11 + 1\cdot 3 + 1\cdot 2 + 1\cdot 6 + 1\cdot 1 = 10\\
C_3(x) & = c_1(x)\cdot 11 + c_3(x)\cdot 3 + c_5(x)\cdot 2 + c_6(x)\cdot 6 + c_7(x)\cdot 1 \\
       & = 10\cdot 11 + 1\cdot 3 + 1\cdot 2 + 1\cdot 6 + 1\cdot 1 = 10\\
C_4(x) & = c_1(x)\cdot 11 + c_3(x)\cdot 3 + c_5(x)\cdot 2 + c_6(x)\cdot 6 + c_7(x)\cdot 1 \\
       & = 10\cdot 11 + 1\cdot 3 + 1\cdot 2 + 1\cdot 6 + 1\cdot 1 = 10
\end{align*}

Теперь, чтобы вычислить элементы $\F_{13}$, заданные этими полиномами, проверяющий должен оценить каждый из этих полиномов на всех шести элементах $T=\{1,2,3,4,5,6\}$ из определения \ref{def:3-fac-zk-circuit}. Это означает, что проверяющий должен вычислить коэффициенты $(a_0,a_1,\ldots,a_5)$ полинома $C(x)=a_0+a_1x+a_2x^2+a_3x^3+a_4x^4+a_5x^5$, который удовлетворяет $C(1)=10$, $C(2)=10$, $C(3)=10$, $C(4)=10$, $C(5)=10$ и $C(6)=10$. В этом простом случае получаем константный полином $C(x)=10$. Проверяющий затем может вычислить правую часть парного уравнения:
\begin{align*}
& e(g^{\alpha},h^{\beta}) \cdot \prod_{j=1}^{4} e(g^{C_j(\tau)},h^{j\cdot \gamma}) \cdot e(\mathbb{C},h^{\delta})\\
& = e(38^2,27^5) \cdot e(38^{C_1(4)},27^{1\cdot 3})\cdot e(38^{C_2(4)},27^{2\cdot 3})\cdot e(38^{C_3(4)},27^{3\cdot 3})\cdot e(38^{C_4(4)},27^{4\cdot 3}) \cdot e(27,27^7)\\
& = e(38^2,27^5) \cdot e(38^{10},27^{3})\cdot e(38^{10},27^{6})\cdot e(38^{10},27^{9})\cdot e(38^{10},27^{12}) \cdot e(27,27^7)\\
& = e(38,27)^{2\cdot 5} \cdot e(38,27)^{10\cdot 3}\cdot e(38,27)^{10\cdot 6}\cdot e(38,27)^{10\cdot 9}\cdot e(38,27)^{10\cdot 12} \cdot e(38,27)^{1\cdot 7}\\
& = e(38,27)^{10 + 4 + 8 + 12 + 3 + 7} = e(38,27)^{4} = 9
\end{align*}
Поскольку обе части парного уравнения равны, проверяющий принимает доказательство как допустимое.

\begin{example}[Схема 3-факторизации Groth16, Фаза верификации]
\label{ex:3-fac-groth-16-verify-circom}
Реализацию фазы верификации Groth16 zk-SNARK в реальных приложениях можно наблюдать, рассматривая реализацию задачи 3-факторизации в Circom из ссылки \ref{ex:3-fac-circom}, ассоциированные R1CS из ссылки \ref{ex:3-fac-r1cs-circom}, параметры из примера \ref{ex:3-fac-groth-16-params-circom}, общую опорную строку из примера \ref{ex:3-fac-groth-16-setup-circom} и zk-SNARK из примера \ref{ex:3-fac-groth-16-prover-circom}.

Предполагая, что экземпляр, сгенерированный доказывающим, дан как $I_1=11$ и zk-SNARK сохранён в файле \fname{proof.json}, а экземпляр сохранён в файле \fname{public.json}, проверяющий может проверить zk-SNARK, выполнив следующую команду:
\\
\\
\texttt{:\$ snarkjs groth16 verify verification\_key.json public.json proof.json}
\\
\\
Команда \texttt{verify} используется для проверки zk-SNARK, и первый параметр после \texttt{verify} относится к ключу проверки из общей опорной строки, который был экспортирован в конце фазы установки. Второй параметр — это публичный вход, а третий — zk-SNARK.
\end{example}

\begin{exercise}
\label{ex:baby-jubjub-circom-verify}
Рассмотрите упражнение \ref{ex:baby-jubjub-circom-setupt} и zk-SNARK, сгенерированный в упражнении \ref{ex:baby-jubjub-circom-setupt}. Проверьте zk-SNARK.
\end{exercise}

\subsection{Симуляция доказательства} Чтобы понять, почему схема Groth16 обеспечивает нулевое разглашение, давайте повторно рассмотрим пример задачи 3-факторизации и предположим, что симулятор узнал секрет параметра $\tau=4$, а также параметры $(\alpha,\beta,\gamma,\delta)=(2,5,3,7)$ из симуляционного люка ссылки \ref{def:groth16-trapdoor}. Кроме того, предположим, что симулятор получил экземпляр $I_1=11$.

Чтобы смоделировать zk-SNARK для этого экземпляра, симулятор сначала выбирает три случайных элемента $a$, $b$ и $c$ из $\F_{13}$ и вычисляет
\begin{align*}
[\mathbb{A}]_1 & = g^{a}\\
               & = 38^{2} = 10
\end{align*}

Затем симулятор вычисляет
\begin{align*}
[\mathbb{B}]_2 & = h^{b}\\
               & = 27^{3} = 18
\end{align*}

Чтобы вычислить третью составляющую, симулятор должен вычислить полиномы $P(x)$, $R(x)$, $O(x)$ и $C(x)$, как описано в фазе доказывающего и фазе верификации. В нашем случае мы получаем:
\begin{align*}
P(x) & = 6x^5 + 8x^4 + 8x^3 + 10x^2 + 7x + 11\\
R(x) & = 8x^5 + 2x^4 + 7x^3 + 4x^2 + 9x + 11\\
O(x) & = 6x^5 + 8x^4 + 8x^3 + 10x^2 + 7x + 11\\
C(x) & = 10
\end{align*}

Затем симулятор вычисляет
\begin{align*}
[\mathbb{C}]_1 & = g^{\frac{a\cdot b - (\beta\cdot P(\tau) + \alpha\cdot R(\tau) + O(\tau) + C(\tau))}{\delta}}\\
               & = 38^{\frac{2\cdot 3 - (5\cdot 8 + 2\cdot 8 + 8 + 10)}{7}}\\
               & = 38^{\frac{6 - (40 + 16 + 8 + 10)}{7}}\\
               & = 38^{\frac{6 - 74}{7}} = 38^{\frac{6 - 9}{7}} = 38^{\frac{10}{7}} = 38^{10\cdot 2}\\
               & = 38^{7} = 27
\end{align*}

Здесь мы использовали тот факт, что $74 = 9$ в $\F_{13}$, и поскольку $6-9 = -3 = 10$ в $\F_{13}$, мы получаем $10/7 = 10\cdot 2 = 7$ в $\F_{13}$, поскольку $7^{-1}=2$.

Таким образом, симулятор вычисляет следующий zk-SNARK:
\begin{equation}
S_{<3\_fac,3\_fac.r1cs,11>} = (10; 18; 27)
\end{equation}

Чтобы проверить, что этот zk-SNARK является допустимым, мы вычисляем обе стороны парного уравнения:
\begin{align*}
e(\mathbb{A},\mathbb{B}) & = e(10,18) = e(38^2,27^3) = e(38,27)^{2\cdot 3} = e(38,27)^6 = 9^6 = 9
\end{align*}

и
\begin{align*}
& e(g^{\alpha},h^{\beta}) \cdot \prod_{j=1}^{4} e(g^{C_j(\tau)},h^{j\cdot \gamma}) \cdot e(\mathbb{C},h^{\delta})\\
& = e(38^2,27^5) \cdot e(38^{10},27^{3})\cdot e(38^{10},27^{6})\cdot e(38^{10},27^{9})\cdot e(38^{10},27^{12}) \cdot e(27,27^7)\\
& = e(38,27)^{10 + 4 + 8 + 12 + 3 + 7} = e(38,27)^{4} = 9
\end{align*}

Поскольку обе части равны, проверяющий принимает доказательство. Однако важно отметить, что этот zk-SNARK был сгенерирован без знания свидетеля $W$. Это демонстрирует, что zk-SNARK не раскрывает никакой информации о свидетеле, поскольку любой zk-SNARK, который может быть сгенерирован с использованием свидетеля, также может быть сгенерирован без него.

\begin{example}[Схема 3-факторизации в Circom и Snark.js, Симуляция]
\label{ex:3-fac-groth-16-simulator-circom} На момент написания Snark.js не имеет алгоритма для генерации смоделированных доказательств с использованием заданного симуляционного люка.
\end{example}

%\subsection{Изменяемость доказательства} Как мы видели в предыдущем параграфе, зная некоторый экземпляр $I$, знание симуляционного люка позволяет вычислить допустимый zk-SNARK без знания фактического свидетеля для этого экземпляра. 

%Другой способ построить 

%Чтобы увидеть это, пусть $\pi=(g_1^A, g_1^C, g_2^B)$ — zk-SNARK для некоторого экземпляра $I$ в протоколе Groth16. Тогда для любого ненулевого элемента поля $a\in\F_r$ другое доказательство $\pi'$ задаётся следующим образом:
%\begin{equation}
%\pi=\Big(\Big(g_1^A\Big)^a, \Big(g_1^C\Big)^a, \Big(g_2^B\Big)^{a^2}\Big)
%\end{equation}  
%Чтобы увидеть это, вспомним из XXX\sme{добавить ссылку}, что проверяющий проверяет допустимость любого заданного SNARK, вычисляя уравнение $e(g_1^A, e_2^B) = e(g_1^\alpha,g_2^\beta)\cdot e(g_1^I,g_2^\gamma)\cdot e(g_1^C,g_2^\delta)$. Однако, если это уравнение выполняется для доказательства $\pi$, оно также выполняется для $\pi'$, поскольку
%\begin{align*}
%e(g_1^A, e_2^B) & = e(g_1^\alpha,g_2^\beta)\cdot e(g_1^I,g_2^\gamma)\cdot e(g_1^C,g_2^\delta) & \Leftrightarrow\\
%\Big(e(g_1^A, e_2^B)\Big)^{a^2} &= \Big(e(g_1^\alpha,g_2^\beta)\cdot e(g_1^I,g_2^\gamma)\cdot e(g_1^C,g_2^\delta)\Big)^{a^2} & \Leftrightarrow\\  
%e((g_1^A)^a, (e_2^B)^a) &= \Big(e(g_1^\alpha,g_2^\beta)\cdot e(g_1^I,g_2^\gamma)\cdot e(g_1^C,g_2^\delta)\Big)^a & \Leftrightarrow\\ 
%\end{align*} 






% ========================= Для версии 2


\begin{comment}
% К сожалению, слишком много для первой версии книги. Закончу это
% в другой раз ... С лучшей ручкой и бумажной хэш-функцией.
\begin{example}[Церемония многостороннего вычисления с обменом участников для схемы факторизации] В этом примере мы хотим смоделировать реальную церемонию многостороннего вычисления с обменом участников для нашей схемы факторизации XXX\sme{добавить ссылку}, как объяснено в XXX\sme{добавить ссылку}.

Мы используем нашу хэш-функцию TinyMD5 XXX\sme{добавить ссылку} для хэширования в $\mathbb{G}_2$.


Мы предполагаем, что у нас есть координатор $Alice$ вместе с тремя сторонами $Bob$, $Carol$ и $Dave$, которые хотят внести свою случайность в протокол. Поскольку степень $n$ целевого полинома равна $2$, нам нужно вычислить \concept{общую опорную строку} (Common Reference String)
$$
CRS= \left\{\right\}
$$
Для того чтобы участник $j>0$ на фазе $l$ вычислил доказательство знания XXX\sme{добавить ссылку}, нам нужно определить $transcript_{l,j-1}$ предыдущего раунда. Мы определяем его как sha256 от $MPC_{l,j-1}$. Более точно, мы определяем
$$
transcript_{1,j-1}= 
MD5(
'[s]g_1 [s]g_2 [s^2] g_1 [\alpha]g_1 [\alpha\cdot s]g_1
[\beta]g_1 [\beta]g_2[\beta \cdot s]g_1'
)
$$
Единственное, что на самом деле важно относительно транскрипта, — это то, что это публично доступные данные, которые не доступны никому до тех пор, пока не существуют данные MPC раунда $j-1$ на фазе $l$.

Мы начинаем с первого раунда, обычно называемого ``степенями тау'' (powers of tau) ОБЪЯСНИТЬ ЭТОТ ТЕРМИН...
Вычисление инициализируется с $s=1$, $\alpha=1$, $\beta=1$. Следовательно, вычисление начинается со следующих данных
$$
MPC_{1,0}= \left\{
\begin{array}{lcl}
([s]g_1, [s]g_2) &=& ((13,15),(7v^2,16v^3))\\ 
{}[s^2] g_1 &=& (13,15)\\
{}[\alpha]g_1 &=& (13,15)\\ 
{}[\alpha\cdot s]g_1 &=& (13,15)\\ 
([\beta]g_1,[\beta]g_2) &=& ((13,15),(7v^2,16v^3))\\ 
{}[\beta \cdot s]g_1 &=& (13,15)
\end{array}
\right\}
$$
Тогда 
\begin{multline*}
transcript_{1,0}=\\ 
MD5('(13,15)(7v^2,16v^3)(13,15)(13,15)(13,15)(13,15)(7v^2,16v^3)(13,15)') =\\ f2baea4d3dba5eef5c63bb210920e7d9
\end{multline*}
Мы получаем этот хэш, вычисляя

$printf '\%s' "(13,15)(7v\textasciicircum 2,16v\textasciicircum 3)(13,15)(13,15)(13,15)(13,15)(7v\textasciicircum 2,16v\textasciicircum 3)(13,15)" | md5sum$
% note the actual code is printf '%s' "(13,15)(7v^2,16v^3)(13,15)(13,15)(13,15)(13,15)(7v^2,16v^3)(13,15)" | md5sum

Все согласились, что MPC начинается 21.03.2020, и каждый может внести вклад в течение ровно года до 20.03.2021. 


  
Затем оно продолжается в стиле round robin, начиная с Bob, который получает эти данные в $MPC_{1,0}$, а затем вычисляет свой вклад. Допустим, что $Bob$ честен и купил 
13-гранный кубик (ИЗОБРАЖЕНИЕ 13-ГРАННОГО КУБИКА), чтобы случайным образом найти три секретных значения поля из нашего простого поля $\F_{13}$. Он бросил кубик и получил $\alpha = 4$, $\beta=8$ и $s= 2$. Затем он обновляет $MPC_{1,0}$:  
$$
MPC_{1,1}= \left\{
\begin{array}{lclcl}
([s]g_1, [s]g_2) &=& ([2](13,15),[2](7v^2,16v^3)) &=& ((33,34),(10v^2,28v^3))\\ 
{}[s^2] g_1 &=& [4](13,15)&=& (35,28)\\
{}[\alpha]g_1 &=& [4](13,15)&=& (35,28)\\ 
{}[\alpha\cdot s]g_1 &=& [8](13,15) &=& (26,9)\\ 
([\beta]g_1,[\beta]g_2) &=& ([8](13,15),[8](7v^2,16v^3))&=& ((26,9),(16v^2,15v^3))\\ 
{}[\beta \cdot s]g_1 &=& [3](13,15)&=& (38,15)
\end{array}
\right\}
$$
В дополнение мы вычисляем следующим образом:
$$
POK_{1,1} \left\{
\begin{array}{lcl}
y_{s} &=& POK(2, f2baea4d3dba5eef5c63bb210920e7d9) = ((33,34),(16v^2 , 28v^3))\\
y_{\alpha} &=& POK(4, f2baea4d3dba5eef5c63bb210920e7d9) = ((35,28),(10v^2 , 15v^3))\\ 
y_{\beta} &=& POK(8, f2baea4d3dba5eef5c63bb210920e7d9) = ((26,9),(16v^2 , 28v^3))\\
\end{array}
\right\}
$$
поскольку $[s]g_1 = (33,34)$, $[\alpha] g_1 = (35,28)$ и $[\beta] g_1 = (26,9)$. а также 
\begin{align*}
TinyMD5_{2}('(33,34)f2baea4d3dba5eef5c63bb210920e7d9') =\\ H_2(MD5('(33,34)f2baea4d3dba5eef5c63bb210920e7d9').trunc(3))=\\ H_2(2066b3b6b6d97c46c3ac6ee2ccd23ad9.trunc(3))= H_2(ad9) = \\
H_2(101 011 011 001)=\\
[8\cdot 4^{1}\cdot 5^{0}\cdot 7^{1}](7v^2 , 16v^3)+
[12\cdot 1^{0}\cdot 3^{1}\cdot 8^{1}](42v^2 , 16v^3 )+\\
[2\cdot 3^{0}\cdot 9^{1}\cdot 11^{1}](17v^2 , 15v^3 ) +
[3\cdot 6^{0}\cdot 9^{0}\cdot 10^{1}](10v^2 , 15v^3 ) =\\
[8\cdot 4\cdot 7](7v^2 , 16v^3)+
[12\cdot 3\cdot 8](42v^2 , 16v^3 )+
[2\cdot 9\cdot 11](17v^2 , 15v^3 ) +
[3\cdot 10](10v^2 , 15v^3 ) =\\
[8\cdot 4\cdot 7](7v^2 , 16v^3)+
[12\cdot 3\cdot 8](42v^2 , 16v^3 )+
[2\cdot 9\cdot 11](17v^2 , 15v^3 ) +
[3\cdot 10](10v^2 , 15v^3 ) =\\
[3](7v^2 , 16v^3)+
[2](42v^2 , 16v^3 )+
[3](17v^2 , 15v^3 ) +
[4](10v^2 , 15v^3 )=\\
[3](7v^2 , 16v^3)+
[2*3](7v^2 , 16v^3 )+
[3*7](7v^2 , 16v^3 ) +
[4*11](7v^2 , 16v^3 )=\\
(42v^2 , 16v^3)+
(17v^2 , 28v^3 )+
(16v^2 , 15v^3 ) +
(16v^2 , 28v^3 )=\\
[3](7v^2 , 16v^3)+
[6](7v^2 , 16v^3 )+
[8](7v^2 , 16v^3 ) +
[5](7v^2 , 16v^3 )=\\
[3+6+8+5](7v^2 , 16v^3)=
(37v^2 , 16v^3 )
\end{align*}
Таким образом, мы получаем $[2](37v^2 , 16v^3 )= (16v^2 , 28v^3 )$

===================

\begin{align*}
TinyMD5_{2}('(35,28)f2baea4d3dba5eef5c63bb210920e7d9') =\\ H_2(MD5('(35,28)f2baea4d3dba5eef5c63bb210920e7d9').trunc(3))=\\ H_2(ad54fa3674f6a84fab9208d7a94c9163.trunc(3))= H_2(163) = \\
H_2(000 101 100 011)=\\
[8\cdot 4^{0}\cdot 5^{0}\cdot 7^{0}](7v^2 , 16v^3)+
[12\cdot 1^{1}\cdot 3^{0}\cdot 8^{1}](42v^2 , 16v^3 )+\\
[2\cdot 3^{1}\cdot 9^{0}\cdot 11^{0}](17v^2 , 15v^3 ) +
[3\cdot 6^{0}\cdot 9^{1}\cdot 10^{1}](10v^2 , 15v^3 ) = \\
[8](7v^2 , 16v^3)+
[12\cdot 8](42v^2 , 16v^3 )+
[2\cdot 3](17v^2 , 15v^3 ) +
[3\cdot 9\cdot 10](10v^2 , 15v^3 ) = \\
[8](7v^2 , 16v^3)+
[5](42v^2 , 16v^3 )+
[6](17v^2 , 15v^3 ) +
[10](10v^2 , 15v^3 ) = \\
[8](7v^2 , 16v^3)+
[5*3](7v^2 , 16v^3 )+
[6*7](7v^2 , 16v^3 ) +
[10*11](7v^2 , 16v^3 )=\\
(16v^2 , 15v^3)+
(10v^2 , 28v^3 )+
(42v^2 , 16v^3 ) +
(17v^2 , 28v^3 )=\\
[8](7v^2 , 16v^3)+
[2](7v^2 , 16v^3 )+
[3](7v^2 , 16v^3 ) +
[6](7v^2 , 16v^3 )=\\
[8+2+3+6](7v^2 , 16v^3)=
(17v^2 , 28v^3 )
\end{align*}
Таким образом, мы получаем $[4](17v^2 , 28v^3 )= (10v^2 , 15v^3 )$

\begin{align*}
TinyMD5_{2}('(26,9)f2baea4d3dba5eef5c63bb210920e7d9') =\\ H_2(MD5('(26,9)f2baea4d3dba5eef5c63bb210920e7d9').trunc(3))=\\ H_2(b87b632f7027ad78cadc2452beb30e9a.trunc(3))= H_2(e9a) = \\
H_2(111 010 011 010)=\\
[8\cdot 4^{1}\cdot 5^{1}\cdot 7^{1}](7v^2 , 16v^3)+
[12\cdot 1^{0}\cdot 3^{1}\cdot 8^{0}](42v^2 , 16v^3 )+\\
[2\cdot 3^{0}\cdot 9^{1}\cdot 11^{1}](17v^2 , 15v^3 ) +
[3\cdot 6^{0}\cdot 9^{1}\cdot 10^{0}](10v^2 , 15v^3 )= \\
[8\cdot 4\cdot 5\cdot 7](7v^2 , 16v^3)+
[12\cdot 3](42v^2 , 16v^3 )+
[2\cdot 9\cdot 11](17v^2 , 15v^3 ) +
[3\cdot 9](10v^2 , 15v^3 )= \\
[2](7v^2 , 16v^3)+
[10](42v^2 , 16v^3 )+
[3](17v^2 , 15v^3 ) +
[1](10v^2 , 15v^3 )= \\
[2](7v^2 , 16v^3)+
[10*3](7v^2 , 16v^3 )+
[3*7](7v^2 , 16v^3 ) +
[1*11](7v^2 , 16v^3 )=\\
(10v^2 , 28v^3)+
(37v^2 , 27v^3 )+
(16v^2 , 15v^3 ) +
(10v^2 , 15v^3 )=\\
[2](7v^2 , 16v^3)+
[4](7v^2 , 16v^3 )+
[8](7v^2 , 16v^3 ) +
[11](7v^2 , 16v^3 )=\\
[2+4+8+11](7v^2 , 16v^3)=
(7v^2 , 27v^3 )
\end{align*}
Таким образом, мы получаем $[8](17v^2 , 28v^3 )= (16v^2 , 28v^3 )$

Таким образом, Bob публикует $MPC_{1,1}$, а также $POK_{1,1}$, и после этого наступает очередь Carol. Допустим также, что Carol честна. Carol смотрит на данные Bob и вычисляет транскрипт в соответствии с нашими правилами
\begin{multline*}
transcript_{1,1}=\\ 
MD5('
(33,34)(10v^2,28v^3)(35,28)(35,28)(26,9)(26,9)(16v^2,15v^3)(38,15)') =\\ fe72e18b90014062682a77136944e362
\end{multline*}
Мы получаем этот хэш, вычисляя

$printf '\%s' "(33,34)(10v^2,28v^3)(35,28)(35,28)(26,9)(26,9)(16v^2,15v^3)(38,15)" | md5sum$

Carol затем вычисляет свой вклад. Поскольку она честна, она случайным образом выбирает три секретных значения поля из нашего простого поля $\F_{13}$, используя свой компьютер. Она нашла $\alpha = 3$, $\beta=4$ и $s= 9$ и обновляет $MPC_{1,1}$:  
$$
MPC_{1,2}= \left\{
\begin{array}{lclcl}
([s]g_1, [s]g_2) &=& ([9](33,34),[9](10v^2,28v^3)) &=&  ((26,34),(16v^2,28v^3))\\ 
{}[s^2] g_1 &=& [9\cdot 9](35,28) &=& (13,28)\\
{}[\alpha]g_1 &=& [3](35,28) &=& (13,28) \\ 
{}[\alpha\cdot s]g_1 &=& [3\cdot 9](26,9) &=& (26,9)\\ 
([\beta]g_1,[\beta]g_2) &=& ([4](26,9),[4](16v^2,15v^3)) &=& ((27,34),(17v^2,28v^3))\\ 
{}[\beta \cdot s]g_1 &=& [4\cdot 9](38,15) &=& (35,28)
\end{array}
\right\}
$$
В дополнение мы вычисляем следующим образом:
$$
POK_{1,2} \left\{
\begin{array}{lcl}
y_{s} &=& POK(9, fe72e18b90014062682a77136944e362) = ((35,15),(17v^2 , 28v^3))\\
y_{\alpha} &=& POK(3, fe72e18b90014062682a77136944e362) = ((38,15),(17v^2 , 15v^3 ))\\ 
y_{\beta} &=& POK(4, fe72e18b90014062682a77136944e362) = ((35,28),(42v^2 , 27v^3 ))\\
\end{array}
\right\}
$$

\begin{align*}
TinyMD5_{2}('(35,15)fe72e18b90014062682a77136944e362') =\\ H_2(MD5('(35,15)fe72e18b90014062682a77136944e362').trunc(3))=\\ H_2(115f145ceffdda73e916dc5ba8ae7354.trunc(3))= H_2(354) = \\
H_2(001 101 010 100)=\\
[8\cdot 4^{0}\cdot 5^{0}\cdot 7^{1}](7v^2 , 16v^3)+
[12\cdot 1^{1}\cdot 3^{0}\cdot 8^{1}](42v^2 , 16v^3 )+\\
[2\cdot 3^{0}\cdot 9^{1}\cdot 11^{0}](17v^2 , 15v^3 ) +
[3\cdot 6^{1}\cdot 9^{0}\cdot 10^{0}](10v^2 , 15v^3 )= \\
[8\cdot 7](7v^2 , 16v^3)+
[12\cdot 8](42v^2 , 16v^3 )+
[2\cdot 9](17v^2 , 15v^3 ) +
[3\cdot 6](10v^2 , 15v^3 )= \\
[4](7v^2 , 16v^3)+
[5](42v^2 , 16v^3 )+
[5](17v^2 , 15v^3 ) +
[5](10v^2 , 15v^3 )= \\
[4](7v^2 , 16v^3)+
[5*3](7v^2 , 16v^3 )+
[5*7](7v^2 , 16v^3 ) +
[5*11](7v^2 , 16v^3 )=\\
(37v^2 , 27v^3)+
(10v^2 , 28v^3 )+
(37v^2 , 16v^3 ) +
(42v^2 , 16v^3 )=\\
[4](7v^2 , 16v^3)+
[2](7v^2 , 16v^3 )+
[9](7v^2 , 16v^3 ) +
[3](7v^2 , 16v^3 )=\\
[4+2+9+3](7v^2 , 16v^3)=
(16v^2 , 28v^3 )
\end{align*}
Таким образом, мы получаем $[9](16v^2 , 28v^3 )= (17v^2 , 28v^3 )$

\begin{align*}
TinyMD5_{2}('(38,15)fe72e18b90014062682a77136944e362') =\\ H_2(MD5('(38,15)fe72e18b90014062682a77136944e362').trunc(3))=\\ H_2(cc4da0c02c4c1b15e72d6cc6430206ab.trunc(3))= H_2(6ab) = \\
H_2(011 010 101 011)=\\
[8\cdot 4^{0}\cdot 5^{1}\cdot 7^{1}](7v^2 , 16v^3)+
[12\cdot 1^{0}\cdot 3^{1}\cdot 8^{0}](42v^2 , 16v^3 )+\\
[2\cdot 3^{1}\cdot 9^{0}\cdot 11^{1}](17v^2 , 15v^3 ) +
[3\cdot 6^{0}\cdot 9^{1}\cdot 10^{1}](10v^2 , 15v^3 )= \\
[8\cdot 5\cdot 7](7v^2 , 16v^3)+
[12\cdot 3](42v^2 , 16v^3 )+
[2\cdot 3\cdot 11](17v^2 , 15v^3 ) +
[3\cdot 9\cdot 10](10v^2 , 15v^3 )= \\
[7](7v^2 , 16v^3)+
[10](42v^2 , 16v^3 )+
[1](17v^2 , 15v^3 ) +
[10](10v^2 , 15v^3 )= \\
[7](7v^2 , 16v^3)+
[10*3](7v^2 , 16v^3 )+
[1*7](7v^2 , 16v^3 ) +
[10*11](7v^2 , 16v^3 )=\\
(17v^2 , 15v^3)+
(17v^2 , 28v^3 )+
(17v^2 , 15v^3 ) +
(17v^2 , 28v^3 )=\\
[7](7v^2 , 16v^3)+
[4](7v^2 , 16v^3 )+
[7](7v^2 , 16v^3 ) +
[6](7v^2 , 16v^3 )=\\
[7+4+7+6](7v^2 , 16v^3)=
(10v^2 , 15v^3)
\end{align*}
Таким образом, мы получаем $[3](10v^2 , 15v^3 )= (17v^2 , 15v^3 )$

\begin{align*}
TinyMD5_{2}('(35,28)fe72e18b90014062682a77136944e362') =\\ H_2(MD5('(35,28)fe72e18b90014062682a77136944e362').trunc(3))=\\ H_2(502323bc55c75f7189fad7999c9f1708.trunc(3))= H_2(708) = \\
H_2(011 100 001 000)=\\
[8\cdot 4^{0}\cdot 5^{1}\cdot 7^{1}](7v^2 , 16v^3)+
[12\cdot 1^{1}\cdot 3^{0}\cdot 8^{0}](42v^2 , 16v^3 )+\\
[2\cdot 3^{0}\cdot 9^{0}\cdot 11^{1}](17v^2 , 15v^3 ) +
[3\cdot 6^{0}\cdot 9^{0}\cdot 10^{0}](10v^2 , 15v^3 )= \\
[8\cdot 5\cdot 7](7v^2 , 16v^3)+
[12](42v^2 , 16v^3 )+
[2\cdot 11](17v^2 , 15v^3 ) +
[3](10v^2 , 15v^3 )= \\
[7](7v^2 , 16v^3)+
[12](42v^2 , 16v^3 )+
[9](17v^2 , 15v^3 ) +
[3](10v^2 , 15v^3 )= \\
[7](7v^2 , 16v^3)+
[12*3](7v^2 , 16v^3 )+
[9*7](7v^2 , 16v^3 ) +
[3*11](7v^2 , 16v^3 )=\\
(17v^2 , 15v^3)+
(42v^2 , 27v^3 )+
(10v^2 , 15v^3 ) +
(17v^2 , 15v^3 )=\\
[7](7v^2 , 16v^3)+
[10](7v^2 , 16v^3 )+
[11](7v^2 , 16v^3 ) +
[7](7v^2 , 16v^3 )=\\
[7+10+11+7](7v^2 , 16v^3)=
(37v^2 , 16v^3)
\end{align*}
Таким образом, мы получаем $[4](37v^2 , 16v^3 )= (42v^2 , 27v^3 )$

Dave думает, что может обхитрить систему. Поскольку он последний, кто вносит вклад, он просто придумывает совершенно новый $MPC$, который не содержит никакой случайности от предыдущих участников. Он думает, что может сделать это, потому что никто не может отличить его $MPC_{1,3}$ от правильного. Если это сделано разумно, он даже сможет вычислить правильные $POK$.

Таким образом, Dave выбирает $s=12$, $\alpha=11$ и $\beta=10$ и сохранит эти значения, надеясь использовать их позже для подделки ложных доказательств в схеме факторизации. Затем он вычисляет  
$$
MPC_{1,3}= \left\{
\begin{array}{lcl}
([s]g_1, [s]g_2) &=& ((13,28),(7v^2,27v^3))\\ 
{}[s^2] g_1 &=& (13,15)\\
{}[\alpha]g_1 &=& (33,9)\\ 
{}[\alpha\cdot s]g_1 &=& (33,34)\\ 
([\beta]g_1,[\beta]g_2) &=& ((38,28),(42v^2,27v^3))\\ 
{}[\beta \cdot s]g_1 &=& (38,15)
\end{array}
\right\}
$$
Dave не удаляет $s$, $\alpha$ и $\beta$, потому что если это будет принято как фаза один вычисления \concept{общей опорной строки} (Common Reference String), Dave уже контролирует 3/4 мошеннического ключа для подделки доказательств. Таким образом, Dave тщательно следит за тем, чтобы доказательства знания были правильными. Он вычисляет транскрипт вклада Carol как 

\begin{multline*}
transcript_{1,2}=\\ 
MD5('
(26,34)(16v^2,28v^3)(13,28)(13,28)(26,9)(27,34)(17v^2,28v^3)(35,28)') =\\ c8e6308fffd47009f5f65e773ae4b499
\end{multline*}

Мы получаем этот хэш, вычисляя

$printf '\%s' "(26,34)(16v^2,28v^3)(13,28)(13,28)(26,9)(27,34)(17v^2,28v^3)(35,28)" | md5sum$

\end{example}
\end{comment}
